\documentclass[12pt,a4paper]{article}

\usepackage[utf8]{inputenc}
\usepackage[T1]{fontenc}
\usepackage[brazilian]{babel}
\usepackage{lmodern}
\usepackage{geometry}
\usepackage{graphicx}
\usepackage{booktabs}
\usepackage{amsmath,amssymb}
\usepackage{hyperref}

\geometry{margin=2.5cm}
\hypersetup{
  colorlinks=true,
  linkcolor=blue,
  citecolor=blue,
  urlcolor=blue,
  pdftitle={Título do Artigo}
}

\DeclareMathOperator{\sen}{sen}

\title{Anotações}
\author{Diego Ferruzzo\\
  \small UFABC \\
  \small \texttt{diego.ferruzzo@ufabc.edu.br}}
\date{\today}

\begin{document}

\maketitle

\begin{abstract}
Anotações sobre o modelo bo Boeing 747.
\end{abstract}

\noindent\textbf{Palavras-chave:} Boeing 747; modelo não linear; Simulação.

\section{Introdução}
Apresente o contexto, a motivação e o objetivo do trabalho. Explique também a
organização das seções seguintes.

\section{Modelo do Boeing 747}
Segundo \cite{franklin2013,brysonControlSpacecraftAircraft1994},
Reescrevemos as equações de movimento como,
\begin{align}
    \begin{split}
        m(\dot{U} + qW - rV) &= X - mg\sen\theta + \kappa T \cos\theta,\\
        m(\dot{V} + rU - pW) &= Y + mg cos\theta \sen\phi,\\
        m(\dot{W} + pV - qU) &= Z + mg cos\theta \cos\phi - \kappa T \sen\theta,\\
        I_x \dot{p} + I_{xz} \dot{r} + (I_z - I_y)qr - I_{xz}qp &= L,\\
        I_y \dot{q} + (I_x - I_z)pr + I_{xz}(r^2 - p^2) &= M,\\
        I_z \dot{r} + I_{xz} \dot{p} + (I_y - I_x)qp - I_{xz}qr &= N,
    \end{split}
    \label{eq:boeing747-nonlinear-model_1}
\end{align}
Reescrevemos as equações de movimento \eqref{eq:boeing747-nonlinear-model_1} como,
\begin{align}
    \begin{split}
        m(\dot{U} + qW - rV) &= X - mg\sen\theta + \kappa T \cos\theta,\\
        \dot{U} + qW - rV &= X/m - g\sen\theta + (1/m)\kappa T \cos\theta,\\
        \dot{U} &= - g\sen\theta + (1/m)\kappa T \cos\theta + X/m - qW + rV ,\\
        %
        m(\dot{V} + rU - pW) &= Y + mg cos\theta \sen\phi,\\
        \dot{V} + rU - pW &= Y/m + g cos\theta \sen\phi,\\
        \dot{V} &= g cos\theta \sen\phi + Y/m - rU + pW ,\\
        %
        m(\dot{W} + pV - qU) &= Z + mg cos\theta \cos\phi - \kappa T \sen\theta,\\
        \dot{W} + pV - qU &= Z/m + g cos\theta \cos\phi - (1/m)\kappa T \sen\theta,\\
        \dot{W} &= g cos\theta \cos\phi - (1/m)\kappa T \sen\theta + Z/m  - pV + qU ,\\
        %
        I_x \dot{p} + I_{xz} \dot{r} + (I_z - I_y)qr - I_{xz}qp &= L,\\
        %
        I_y \dot{q} + (I_x - I_z)pr + I_{xz}(r^2 - p^2) &= M,\\
        I_y \dot{q} &= M - (I_x - I_z)pr - I_{xz}(r^2 - p^2) ,\\
        \dot{q} &= M/I_y  - ((I_x - I_z)/I_y) pr - (I_{xz}/I_y) (r^2 - p^2) ,\\
        %
        I_z \dot{r} + I_{xz} \dot{p} + (I_y - I_x)qp - I_{xz}qr &= N,
    \end{split}
    \label{eq:boeing747-nonlinear-model_2}
\end{align}
das duas equações a seguir determinamos $\dot{p}$ e $\dot{r}$,
\begin{align}
    \begin{split}
       I_x \dot{p} + I_{xz} \dot{r} + (I_z - I_y)qr + I_{xz}qp &= L,\\
       I_z \dot{r} + I_{xz} \dot{p} + (I_y - I_x)qp - I_{xz}qr &= N,\\
    \end{split}
    \label{eq:pr_acopladas}
\end{align}
Multiplicando a primeira equação por $I_z$ e a segunda por $I_{xz}$, temos:
\begin{align*}
    \begin{split}
        I_zI_x \dot{p} + I_zI_{xz} \dot{r} + I_z(I_z - I_y)qr + I_zI_{xz}qp &= I_zL,\\
        I_{xz}I_z \dot{r} + I_{xz}^2 \dot{p} + I_{xz}(I_y - I_x)qp - I_{xz}^2qr &=I_{xz}N,\\
    \end{split}
\end{align*}
Subtraindo a segunda equação da primeira, temos:
\begin{align*}
         (I_zI_x - I_{xz}^2) \dot{p} + (I_z(I_z - I_y) + I_{xz}^2)qr + (I_zI_{xz} - I_{xz}(I_y - I_x))qp &= I_zL - I_{xz}N,
\end{align*}
Isolando $\dot{p}$, temos:
\begin{align*}
         (I_zI_x - I_{xz}^2) \dot{p} &= I_zL - I_{xz}N - (I_z(I_z - I_y) + I_{xz}^2)qr - (I_zI_{xz} - I_{xz}(I_y - I_x))qp,
\end{align*}
\begin{align*}
         \dot{p} &=\left(\dfrac{1}{I_zI_x - I_{xz}^2}\right)\left( I_zL - I_{xz}N - (I_z(I_z - I_y) + I_{xz}^2)qr - (I_zI_{xz} - I_{xz}(I_y - I_x))qp\right),
\end{align*}
\begin{align*}
    \boxed{\dot{p} =\left(\dfrac{1}{I_zI_x - I_{xz}^2}\right)\left( I_zL - I_{xz}N - (I_z(I_z - I_y) + I_{xz}^2)qr - I_{xz}(I_z -I_y + I_x)qp\right)},
\end{align*}
Para obter \(\dot{r}\), multiplicamos a primeira em equação, \eqref{eq:pr_acopladas}, por \(I_{xz}\) e a segunda por \(I_x\):
\begin{align*}
    \begin{split}
       I_{xz}I_x \dot{p} + I_{xz}^2 \dot{r} + I_{xz}(I_z - I_y)qr + I_{xz}^2 qp &= I_{xz} L,\\
       I_xI_z \dot{r} + I_xI_{xz} \dot{p} + I_x(I_y - I_x) qp - I_xI_{xz} qr &= I_x N,
    \end{split}
\end{align*}
Substituindo a segunda equação na primeira, temos:
\begin{align*}
        (I_{xz}^2-I_xI_z)\dot{r} + (I_{xz}(I_z - I_y) +  I_xI_{xz}) qr + (I_{xz}^2  - I_x(I_y - I_x)) qp &= I_{xz} L - I_x N,
   \end{align*}
\begin{align*}
        (I_{xz}^2-I_xI_z)\dot{r} &= I_{xz} L - I_x N - (I_{xz}(I_z - I_y) +  I_xI_{xz}) qr - (I_{xz}^2  - I_x(I_y - I_x)) qp ,
   \end{align*}
\begin{align*}
    \boxed{\dot{r} = \left(\dfrac{1}{I_{xz}^2-I_xI_z}\right)\left(I_{xz} L - I_x N - (I_{xz}(I_z - I_y) +  I_xI_{xz}) qr - (I_{xz}^2  - I_x(I_y - I_x)) qp\right),}
   \end{align*}
Consequentemente, as equações de movimento podem ser reescritas como:
\begin{align}
    \begin{split}
        \dot{U} &= - g\sen\theta + (1/m)\kappa T \cos\theta - qW + rV + X/m  ,\\
        \dot{V} &= g cos\theta \sen\phi  - rU + pW + Y/m\\
        \dot{W} &= g cos\theta \cos\phi - (1/m)\kappa T \sen\theta  - pV + qU + Z/m ,\\
        \dot{p} & =\left(\dfrac{1}{I_zI_x - I_{xz}^2}\right)\left( I_zL - I_{xz}N - (I_z(I_z - I_y) + I_{xz}^2)qr - I_{xz}(I_z -I_y + I_x)qp\right),\\
        \dot{q} &= M/I_y  - ((I_x - I_z)/I_y) pr - (I_{xz}/I_y) (r^2 - p^2) ,\\
        \dot{r} &= \left(\dfrac{1}{I_{xz}^2-I_xI_z}\right)\left(I_{xz} L - I_x N - (I_{xz}(I_z - I_y) +  I_xI_{xz}) qr - (I_{xz}^2  - I_x(I_y - I_x)) qp\right).
    \end{split}
    \label{eq:boeing747-nonlinear-model_3}
\end{align}
Segundo Bryson \cite{brysonControlSpacecraftAircraft1994}, a relação entre as variações dos ângulos de Euler e as velocidades angulares do corpo é dada por:
\begin{align}
    \begin{split}
        \dot{\phi} &= p + q\sen\phi\tan\theta + r\cos\phi\tan\theta,\\
        \dot{\theta} &= q\cos\phi - r\sen\phi,\\
        \dot{\psi} &= q\sen\phi/\cos\theta + r\cos\phi/\cos\theta.
    \end{split}
    \label{eq:boeing747-euler-angles}
\end{align}
\bibliographystyle{plain}
\bibliography{referencias}

\end{document}
